% Propietario conservado: /Users/ruben/Documents/New project/output/REVISION_ARGUMENTAL_APERTURAS_CIERRES_20260915/VI/delivery/MANUSCRIPT_VI_ES_EN_COMPLETE/source_es/sections/14_apendice_censo_regional.tex
\section{Censo completo de firmas y órbitas regionales}
\label{vi:vi:reg:tablas}

Las tablas siguientes completan la prueba finita de la sección
\ref{vi:vi:reg:seccion}. Cada firma se obtiene de
\eqref{vi:vi:reg:firmas}; su multiplicidad cuenta las condiciones
iniciales de un cursor, no parejas de cursores. Las parejas se
forman después mediante \eqref{vi:eq:emisor-seis}.

\subsection{Las firmas de las dos hojas}

% BEGIN OWNER /Users/ruben/Documents/New project/output/REVISION_ARGUMENTAL_APERTURAS_CIERRES_20260915/VI/delivery/MANUSCRIPT_VI_ES_EN_COMPLETE/source_es/sections/14_tablas_firmas_regionales.tex
% Generado por technical/generar_catalogo_regional.py; censo completo de firmas.
\begin{longtable}{@{}llr@{}}
\toprule Hoja&Firma de seis ventanas&Preimágenes\\\midrule\endhead
Aditiva&\(\mathtt{122564}\)&18\\
Aditiva&\(\mathtt{122898}\)&18\\
Aditiva&\(\mathtt{212645}\)&18\\
Aditiva&\(\mathtt{212988}\)&18\\
Aditiva&\(\mathtt{221456}\)&18\\
Aditiva&\(\mathtt{221889}\)&18\\
Aditiva&\(\mathtt{456221}\)&18\\
Aditiva&\(\mathtt{465898}\)&18\\
Aditiva&\(\mathtt{546988}\)&18\\
Aditiva&\(\mathtt{564122}\)&18\\
Aditiva&\(\mathtt{645212}\)&18\\
Aditiva&\(\mathtt{654889}\)&18\\
Aditiva&\(\mathtt{889221}\)&18\\
Aditiva&\(\mathtt{889654}\)&18\\
Aditiva&\(\mathtt{898122}\)&18\\
Aditiva&\(\mathtt{898465}\)&18\\
Aditiva&\(\mathtt{988212}\)&18\\
Aditiva&\(\mathtt{988546}\)&18\\
Multiplicativa&\(\mathtt{000000}\)&108\\
Multiplicativa&\(\mathtt{177393}\)&8\\
Multiplicativa&\(\mathtt{177555}\)&8\\
Multiplicativa&\(\mathtt{177771}\)&8\\
Multiplicativa&\(\mathtt{339555}\)&8\\
Multiplicativa&\(\mathtt{339771}\)&8\\
Multiplicativa&\(\mathtt{339933}\)&8\\
Multiplicativa&\(\mathtt{393177}\)&8\\
Multiplicativa&\(\mathtt{393393}\)&8\\
Multiplicativa&\(\mathtt{393555}\)&8\\
Multiplicativa&\(\mathtt{555177}\)&8\\
Multiplicativa&\(\mathtt{555339}\)&8\\
Multiplicativa&\(\mathtt{555393}\)&8\\
Multiplicativa&\(\mathtt{555555}\)&24\\
Multiplicativa&\(\mathtt{555717}\)&8\\
Multiplicativa&\(\mathtt{555771}\)&8\\
Multiplicativa&\(\mathtt{555933}\)&8\\
Multiplicativa&\(\mathtt{717555}\)&8\\
Multiplicativa&\(\mathtt{717717}\)&8\\
Multiplicativa&\(\mathtt{717933}\)&8\\
Multiplicativa&\(\mathtt{771177}\)&8\\
Multiplicativa&\(\mathtt{771339}\)&8\\
Multiplicativa&\(\mathtt{771555}\)&8\\
Multiplicativa&\(\mathtt{933339}\)&8\\
Multiplicativa&\(\mathtt{933555}\)&8\\
Multiplicativa&\(\mathtt{933717}\)&8\\
\bottomrule
\end{longtable}

% END OWNER


\subsection{Las órbitas de la imagen ternaria}
El representante es el mínimo lexicográfico de la órbita. «Fibra»
cuenta emisiones decimales sobre una palabra y «Peso» cuenta sus
condiciones iniciales conjuntas. Las diagonales se reúnen sobre
toda la órbita. El censo \(\nu=(n_0,n_1,n_2)\) conserva la
multiplicidad de símbolos y es invariante bajo las permutaciones.

% BEGIN OWNER /Users/ruben/Documents/New project/output/REVISION_ARGUMENTAL_APERTURAS_CIERRES_20260915/VI/delivery/MANUSCRIPT_VI_ES_EN_COMPLETE/source_es/sections/14_tabla_orbitas_regionales.tex
% Generado por technical/generar_catalogo_regional.py; 43 órbitas, sin omisiones.
\begin{longtable}{@{}lrrrrl@{}}
\toprule Representante&Órbita&Fibra&Peso&\(\nu\)&Diagonales\\\midrule\endhead
\(\mathtt{001002}\)&6&2&288&\((4,1,1)\)&\(\{0,1,3,4,6,7\}\)\\
\(\mathtt{001011}\)&6&2&288&\((3,3,0)\)&\(\{0,1,3,4,6,7\}\)\\
\(\mathtt{001012}\)&6&2&288&\((3,2,1)\)&\(\{2,5,8\}\)\\
\(\mathtt{001021}\)&6&2&2088&\((3,2,1)\)&\(\{2,5,8\}\)\\
\(\mathtt{001022}\)&6&1&144&\((3,1,2)\)&\(\{2,5,8\}\)\\
\(\mathtt{001100}\)&3&6&864&\((4,2,0)\)&\(\{2,5,8\}\)\\
\(\mathtt{001101}\)&6&2&288&\((3,3,0)\)&\(\{0,1,3,4,6,7\}\)\\
\(\mathtt{001110}\)&6&1&1944&\((3,3,0)\)&\(\{0,3,6\}\)\\
\(\mathtt{001112}\)&6&5&1008&\((2,3,1)\)&\(\{0,1,2,3,4,5,6,7,8\}\)\\
\(\mathtt{001202}\)&6&1&144&\((3,1,2)\)&\(\{2,5,8\}\)\\
\(\mathtt{001210}\)&6&1&144&\((3,2,1)\)&\(\{2,5,8\}\)\\
\(\mathtt{001211}\)&6&2&576&\((2,3,1)\)&\(\{2,5,8\}\)\\
\(\mathtt{001220}\)&6&6&864&\((3,1,2)\)&\(\{0,1,3,4,6,7\}\)\\
\(\mathtt{001222}\)&6&1&144&\((2,1,3)\)&\(\{2,5,8\}\)\\
\(\mathtt{002012}\)&6&1&144&\((3,1,2)\)&\(\{1,4,7\}\)\\
\(\mathtt{002112}\)&6&1&144&\((2,2,2)\)&\(\{0,3,6\}\)\\
\(\mathtt{002211}\)&6&1&144&\((2,2,2)\)&\(\{1,4,7\}\)\\
\(\mathtt{002220}\)&6&1&144&\((3,0,3)\)&\(\{0,3,6\}\)\\
\(\mathtt{011021}\)&6&2&288&\((2,3,1)\)&\(\{0,1,3,4,6,7\}\)\\
\(\mathtt{011112}\)&6&2&288&\((1,4,1)\)&\(\{0,1,3,4,6,7\}\)\\
\(\mathtt{011120}\)&6&1&144&\((2,3,1)\)&\(\{0,3,6\}\)\\
\(\mathtt{011201}\)&6&1&1944&\((2,3,1)\)&\(\{1,4,7\}\)\\
\(\mathtt{011211}\)&6&1&144&\((1,4,1)\)&\(\{0,3,6\}\)\\
\(\mathtt{011220}\)&6&1&144&\((2,2,2)\)&\(\{1,4,7\}\)\\
\(\mathtt{011222}\)&6&1&144&\((1,2,3)\)&\(\{1,4,7\}\)\\
\(\mathtt{012112}\)&6&3&432&\((1,3,2)\)&\(\{1,2,4,5,7,8\}\)\\
\(\mathtt{012121}\)&6&1&144&\((1,3,2)\)&\(\{0,3,6\}\)\\
\(\mathtt{012202}\)&6&1&144&\((2,1,3)\)&\(\{0,3,6\}\)\\
\(\mathtt{012210}\)&6&1&144&\((2,2,2)\)&\(\{2,5,8\}\)\\
\(\mathtt{012220}\)&6&1&144&\((2,1,3)\)&\(\{1,4,7\}\)\\
\(\mathtt{012222}\)&6&1&144&\((1,1,4)\)&\(\{2,5,8\}\)\\
\(\mathtt{021022}\)&6&1&144&\((2,1,3)\)&\(\{2,5,8\}\)\\
\(\mathtt{021121}\)&6&3&432&\((1,3,2)\)&\(\{0,1,2,3,4,5,6,7,8\}\)\\
\(\mathtt{021202}\)&6&3&432&\((2,1,3)\)&\(\{0,1,2,3,4,5,6,7,8\}\)\\
\(\mathtt{022022}\)&3&2&288&\((2,0,4)\)&\(\{1,4,7\}\)\\
\(\mathtt{022112}\)&6&1&144&\((1,2,3)\)&\(\{1,4,7\}\)\\
\(\mathtt{022121}\)&6&6&864&\((1,2,3)\)&\(\{0,1,2,3,4,5,6,7,8\}\)\\
\(\mathtt{022202}\)&3&2&288&\((2,0,4)\)&\(\{0,3,6\}\)\\
\(\mathtt{022211}\)&6&2&288&\((1,2,3)\)&\(\{1,2,4,5,7,8\}\)\\
\(\mathtt{022222}\)&6&1&144&\((1,0,5)\)&\(\{1,4,7\}\)\\
\(\mathtt{112112}\)&3&6&1440&\((0,4,2)\)&\(\{0,1,2,3,4,5,6,7,8\}\)\\
\(\mathtt{112211}\)&3&4&576&\((0,4,2)\)&\(\{1,2,4,5,7,8\}\)\\
\(\mathtt{112222}\)&6&2&288&\((0,2,4)\)&\(\{1,2,4,5,7,8\}\)\\
\bottomrule
\end{longtable}

% END OWNER

